\documentclass[12pt]{article} 
\renewcommand{\thempfootnote}{\arabic{mpfootnote}}

\usepackage[utf8]{inputenc}

\usepackage{multicol}
\setlength{\columnsep}{1cm}

\usepackage{abc}
\usepackage{graphicx}
\usepackage{musicography}

\usepackage[margin=0.5in]{geometry} 
\geometry{letterpaper}

\usepackage{fancyhdr} 
\pagestyle{fancy} % options: empty , plain , fancy
\renewcommand{\headrulewidth}{0pt} % customise the layout...
\lhead{}\chead{}\rhead{}
\lfoot{}\cfoot{\thepage}\rfoot{}

\begin{document}

\thispagestyle{empty}

\begin{abc}[name=greasycoat,options={-O= +c -F single.fmt}]
X:1
T:Old Greasy Coat
C:from Edden Hammons (1875-1955)
Z:abc-transcription: Josh Larios <jdlarios@gmail.com>, 2021.03.23
M:C|
K:AMix
|: AB | "A"c2 A2 G2 AB | c2 A2 A2 AB | c2 A2 GBAG | "E"E2 A2 "A"A2 AB | 
c2 A2 GABd | "D"e2 ef e2 d2 | "A"c2 A2 "G"EGAG | "E"E2 A2 "A"A2 :|: A2- |
"A"ABcd edef | "G"g2 gf g2 A2- | ABcd ede^g | "A"a2 (3ba^g a2 A2- | 
ABcd e2 e2 | "G"a^ga^g e2 d2 | "A"c2 A2 "G"GBAG | "E"E2 A2 "A"A2 :|
\end{abc}
% \vskip 0.5in
\noindent\makebox[\textwidth][c]{%
\begin{minipage}{7.0in}
  \setlength{\parindent}{1.5em}

  This version of ``Greasy Coat'' comes from the playing of Edden
  Hammons, of West Virginia.  It's played in a ``modal'' style with
  the third of the scale often played a little flat (somewhere between
  C\musSharp{} and C\musNatural{} in this case) and the seventh
  sometimes but not always flat (G\musSharp{} or G\musNatural{}). Accompanists might want to choose
  chord shapes which omit the third entirely. If you listen to the
  1947 source recording\footnote{https://www.slippery-hill.com/content/old-greasy-coat}
  you can hear trills in Hammons' playing, showing a Scottish influence.

  There has been somewhat heated debate around the meaning of the term
  ``greasy coat'' over the years. A common claim is that it is an
  old-time euphemism for a condom, although nothing in the written
  record that we've been able to find backs that up. From a search of
  historical books and newspapers, it seems more likely to us that a
  ``greasy coat'' was commonly used as metonymy for someone
  low-class: a hobo or manual laborer with a worn, dirty coat.
  
  If it was ever in common use as a euphemism for a condom, it may
  have been a specific regional term from in or around West
  Virginia. According to folklorist, musicologist, and musician Gerry
  Milnes:
  \begin{quote}
    There was an old banjo player in Greenbrier Co, WV who played it,
    Lee Blankenship. Once, about 1970, after playing he uttered
    something like, ``You don't want to mess with those rubber
    things.'' ``Rubbers'' being the operative word for condoms then.
  \end{quote}
  Greenbrier County is just south of the area where the Hammons family lived.

  Burl Hammons, a nephew of Edden, was recorded around twenty years
  later performing a somewhat different version of ``Greasy Coat'',
  which you can listen to at slippery-hill.com. Lyrics were shared
  with collectors by the Hammons family, but not recorded by them:

  \begin{quote}
    Well I don't drink and I don't smoke,\\
    And I don't mess with the greasy coat.\\
    I don't spit and I don't chew,\\
    And I don't go with the girls that do.\\
    I don't kiss and I don't tell,\\
    And all you sinners gonna go to hell.
  \end{quote}
\end{minipage}
}


\end{document}
